% Tree of Knowledge v5 --- Forest of Knowledge: variable-depth forking with learntropy-driven adaptation
% Core claim: different knowledge types require different fork depths; learntropy drives all decisions.

\documentclass{article}
\usepackage[preprint,nonatbib]{neurips_2024}
\usepackage[utf8]{inputenc}
\usepackage[T1]{fontenc}
\usepackage{lmodern}
\usepackage{hyperref}
\usepackage{url}
\usepackage{booktabs}
\usepackage{amsfonts}
\usepackage{amsmath}
\usepackage{nicefrac}
\usepackage{microtype}
\usepackage{graphicx}
\usepackage{xcolor}
\usepackage{enumitem}
\usepackage{tikz}
\usetikzlibrary{arrows.meta}
\usepackage{wrapfig}

\newcommand{\placeholder}[1]{\begin{center}\fbox{\parbox{0.85\textwidth}{\centering\textcolor{red}{\textbf{[FIGURE: #1]}}}}\end{center}}
\newcommand{\tok}{Forest of Knowledge}

\title{Forest of Knowledge: Learntropy-Driven Separation of \\ Fluid and Crystallized Intelligence in Language Models}

\author{
  Erik de Bruijn \\
  Independent Researcher \\
  \texttt{erik@erikdebruijn.nl}
}

\begin{document}
\maketitle

% ============================================================================
\begin{abstract}
Language models entangle reasoning ability (\emph{fluid intelligence}) with
factual knowledge (\emph{crystallized intelligence}) in shared parameters.
We separate them using variable-rank LoRA expert adapters on FFN layers,
routed by Gumbel-softmax, and propose \emph{learntropy}---per-token
cross-entropy weighted by routing contribution---as the unified control
signal for \emph{when} to fork (bimodal learntropy), \emph{where} to fork
(per-layer divergence), and \emph{how fast} to learn (per-expert learning
rate modulation).

Across 50+ experiments on Qwen3-1.7B and Qwen3-8B, we establish six
findings. First, \emph{causal locality}---selective ablation damage, not
weight orthogonality---is the correct metric for modularity, and requires
sufficient model scale ($M_{ij}$ CV 0.16 at 1.7B vs.\ 0.50--0.67 at 8B).
Second, learntropy-modulated per-expert learning rates break the
dead-expert problem through a \emph{Piagetian inversion}: slow-learning
experts crystallize into the most critical components while
over-accommodating experts go dormant. Third, domain selectivity is
\emph{non-monotonic} in adapter rank, peaking at rank-64 (1.6\% of hidden
dimension) and declining at higher ranks as overcapacity eliminates
specialization pressure. Fourth, \emph{learntropy-directed scheduled
splitting}---splitting the highest-learntropy expert every $N$ steps---grows
a tree from 2 to 8 experts with all 8 active ($M_{ij}$ CV up to 2.34,
reproduced), reducing perplexity by 22.8\% compared to the base model.
Fifth, when combined with learntropy-driven rank increases, the system
autonomously discovers the optimal rank (4$\to$64) through 4 doublings,
converging to the empirical peak of the rank-selectivity curve without
explicit guidance. Sixth, a per-layer \emph{delta gate} on single domain
adapters produces robust domain selectivity ($+0.621 \pm 0.009$, 4~seeds)
while improving generic perplexity by 6\%, enabling hot-pluggable deployment.

The framework separates fluid intelligence (shared trunk) from crystallized
knowledge (expert overrides) through an inheritance mechanism: each expert
inherits the full base model and applies selective LoRA modifications at
variable depth and rank, driven entirely by learntropy signals.
\end{abstract}

% ============================================================================
\section{Introduction}
\label{sec:intro}

Current language models embed all knowledge---from common syntax to rare medical terminology---in the same parameters that perform reasoning. Adding knowledge means adding parameters, and therefore inference cost. The standard alternatives do not escape this trade-off: retrieval-augmented generation adds context tokens, tool calls add generation tokens, and longer contexts scale quadratically in attention. Every path to more knowledge costs inference budget.

The \tok{} framework aims to weaken this coupling by placing domain knowledge in separately-loadable adapters whose cost is zero until activated. The dense core handles reasoning (fluid intelligence) and is always resident in VRAM. Crystallized knowledge lives in lightweight, hot-loadable LoRA adapters. Adding specialists does not slow the core's inference.

A key insight from our experiments is that not all knowledge types diverge at the same depth. Code diverges from natural language in early layers (different syntax, operators, control flow), while medical and legal text share the reasoning trunk through layer~20+ and diverge only in domain-specific vocabulary. This motivates a \emph{forest} rather than a single tree: each knowledge type grows its own branching structure with a fork depth determined by how early its representations diverge from the shared trunk. A per-token difficulty signal---\emph{learntropy}---drives all adaptation decisions: when to fork (bimodal distribution), where to fork (per-layer divergence), and how fast each expert learns (learning rate modulation).

Prior work shows that learned models reside on a low intrinsic
dimension~\cite{li2018intrinsic,aghajanyan2021intrinsic}. We observe that
specialization depth, adapter rank, and storage tier are not three independent
design decisions---they correlate strongly. We distinguish two regimes:

\begin{itemize}[nosep,leftmargin=1em]
\item \textbf{Lens-like} (low rank, e.g.\ rank-1 to rank-4): the adapter does not alter the computation itself but \emph{redirects information} through the existing weights. The original FFN does the heavy lifting; the LoRA adapter bends the light. At 4--16\,KB, these corrections are always cached and cost nothing to load.
\item \textbf{Crystal-like} (high rank, e.g.\ rank-32 to rank-64): the $A \times B$ matrix is large enough to encode its own features---knowledge the original FFN does not possess. It is no longer a correction but an independent information carrier. At 128--256\,KB, these specialists are loaded from flash on demand.
\end{itemize}

The rank is not a hyperparameter to tune---it is a \emph{measurement} of the
dimensionality of the update-subspace needed beyond the parent representation.
A low rank means residual corrections fit in a small subspace (assimilation in
Piaget's terms~\cite{piaget1952}); a high rank means the expert requires
structurally new capacity that the parent cannot provide (accommodation).

\input{figure1.tex}

Standard Mixture-of-Experts (MoE) architectures attempt this separation but
fail in practice: load-balancing losses enforce near-uniform routing, preventing
any hot/cold distinction, and copy-initialized experts remain virtually
identical. We dedicate Section~\ref{sec:why-moe-fails} to this failure mode,
which motivates every design choice in \tok{}.

\paragraph{Concrete benefits.}
\begin{itemize}
    \item Specialists measured in \textbf{kilobytes}, not megabytes (4\,KB--256\,KB vs.\ 16\,MB per MoE expert).
    \item Trainable on \textbf{commodity hardware}, in parallel---each adapter trains independently.
    \item \textbf{Composable by a community}---contributors add specialists without retraining the core.
    \item \textbf{No inference-latency overhead}---the core runs at full speed; adapters merge into the forward pass.
\end{itemize}

\paragraph{Distributed training at community scale.}
Because each LoRA adapter after the initial fork is completely independent,
training can be fully distributed. Individual contributors or volunteer
nodes only need the frozen trunk and a single target adapter (a few
kilobytes to a few hundred kilobytes). They train their assigned specialist
on ZPD-guided, topic-partitioned data using only learntropy-driven signals
and contrastive loss, then return the updated adapter. No synchronization
between experts is required except for occasional lightweight router
updates. This SETI@home-style paradigm makes it possible for a global
community to grow the forest incrementally on commodity hardware, turning
specialization into a collaborative, massively parallel effort.

Sections~\ref{sec:why-moe-fails}--\ref{sec:analysis} use \textbf{Qwen3-1.7B}
(28 layers, 2.0B parameters) as the running example, training on C4 data.
Section~\ref{sec:8b} presents the primary results on \textbf{Qwen3-8B} (36
layers), where causal locality is achieved. All reported numbers are from
Qwen3-1.7B unless otherwise noted.

% ============================================================================
\section{Why Can't Standard MoE Differentiate?}
\label{sec:why-moe-fails}

The promise of Mixture-of-Experts is that different experts will learn different
things. In practice, they do---but not in the way we expect. Experts readily
specialize along \emph{functional} axes (structure vs.\ content, punctuation
vs.\ rare words) but resist specializing along \emph{domain} axes (medical
vs.\ legal vs.\ code). The word ``expert'' evokes a human domain specialist,
but MoE experts behave more like specialized tools: one handles
fasteners, another handles cutting, regardless of whether the project is
carpentry or plumbing. We spent six experimental arms and 40+ hours of GPU
training on Qwen3-1.7B discovering why domain specialization is harder than
functional specialization, and what it takes to cross that boundary.

\paragraph{The copy-initialization trap.}
The standard upcycling recipe~\cite{komatsuzaki2023sparse} creates experts by
copying the dense model's FFN weights. Each expert starts in the same basin of
attraction. Our decisive experiment: we forced completely disjoint data
partitions across experts, training for 12M tokens. Result: expert weight cosine
similarity remained at \textbf{0.998}. The copy initialization creates a basin
that disjoint data alone cannot escape on practical timescales.

\paragraph{Load-balancing destroys structure.}
MoE training adds an auxiliary load-balancing loss that penalizes uneven routing.
This is necessary to prevent expert collapse, but it has a side effect: it
enforces near-uniform token distribution across experts, eliminating the very
frequency structure that would allow hot/cold tiering. Every expert sees the
same distribution and converges to the same function.

\paragraph{Six arms, one conclusion.}
Table~\ref{tab:lessons} summarizes six experimental arms that systematically
explored alternatives. Prescriptive cost losses create routing structure (Gini
0.60) but the effect dilutes over training. Descriptive losses fail entirely
(Gini 0.05). Biological mechanisms help marginally (Gini 0.46). Without
explicit divergence pressure on the \emph{weights themselves}, no genuine
specialization emerges.

\begin{table}[t]
\centering
\caption{Six experimental arms on Qwen3-1.7B that establish what does \emph{not}
produce genuine expert differentiation.}
\label{tab:lessons}
\begin{tabular}{@{}llp{6.5cm}@{}}
\toprule
Arm & Key metric & Lesson \\
\midrule
A: Prescriptive cost loss  & Gini 0.60$\to$0.49 & Creates routing skew, but skew dilutes over training \\
B: Reuse-distance proxy     & Gini 0.05           & Descriptive loss produces no structure \\
C: Adaptive + cost          & 25 rounds, Gini dilutes & Piaget mechanism works, but routing Gini is wrong metric \\
D: Pure data-driven         & Gini 0.11           & No specialization without resource constraint \\
E: Bio-integrated           & Gini 0.46           & Six bio mechanisms help modestly, not fundamentally \\
F: KD-Warm upcycle          & CosSim 0.998        & Copy-init experts don't differentiate via data alone \\
\bottomrule
\end{tabular}
\end{table}

These failures taught us that expert differentiation requires three forces acting
\emph{together}: (1)~different starting points (not copy-initialization),
(2)~different training data, and (3)~explicit divergence pressure on the weights.
The \tok{} framework instantiates all three through learntropy-driven progressive
splitting with contrastive LoRA adapters.

\paragraph{Three levels of differentiation.}
Our experiments reveal that ``expert specialization'' is not binary but admits
at least three distinct levels, each progressively harder to achieve:

\begin{enumerate}[nosep]
    \item \textbf{Parameter differentiation.} Expert weight matrices become
          orthogonal in weight space (low CosSim). Achievable via contrastive
          loss. \emph{Necessary but not sufficient.}
    \item \textbf{Routing differentiation.} The router sends different token
          populations to different experts based on meaningful features (not
          just load balancing). Requires informative hidden states at the
          routing point (Section~\ref{sec:divergence}).
    \item \textbf{Causal modularity.} Each expert must be selectively
          indispensable for a specific token population: removing it causes
          concentrated damage to that population while leaving others
          unaffected. This is the requirement for hot-pluggability, and it is
          fundamentally different from generic otherness---experts must not
          merely \emph{differ}, they must \emph{own} distinct subsets of the
          model's competence. Measured by the ablation matrix $M_{ij}$, not by
          weight geometry.
\end{enumerate}

Our framework achieves level~1 reliably (CosSim $<$0.3) at 1.7B scale,
and level~2 (routing selectivity) via Gumbel-softmax routing. At 8B scale,
level~3 (causal locality) is achieved: $M_{ij}$ CV reaches 0.50--0.67 with
ablation damage varying 3.5--13.5$\times$ across domains
(Section~\ref{sec:8b}). The progression across levels and scales is a
central finding of this paper.

\paragraph{Domains may not be the natural primitive.}
A recurring theme across our experiments and the literature~\cite{standing2026,
domaindriver2026, openmoe2024} is that functional axes (syntax, token type,
sentence position) are the \emph{primary} basis of expert specialization, while
domain axes (medical, legal, code) are secondary and may not emerge as a
natural partition at all. If this is correct, the right framing is not ``how
do we make experts specialize by domain'' but ``what are the natural factors
along which expert specialization occurs, and are domains a meaningful
downstream consequence of those factors?''

\paragraph{Unifying efficiency mechanisms in a single structure.}
Several efficiency mechanisms that MoE systems often introduce as separate
architectural components---such as shared always-on
experts~\cite{deepseek2024moe}, driver-like experts with disproportionate
causal influence~\cite{domaindriver2026}, and selectively routed
specialists---may instead correspond to different depths in a single tree. The
trunk provides a shared substrate; early, low-rank forks can serve as
lightweight infrastructure; and deeper, higher-rank branches can implement
conditional residual specialization. The 8B layer divergence profile (Section~\ref{sec:8b}) provides initial
evidence: the U-shaped fisher ratio pattern (high at layers 0--4 and 24+,
low in the middle) suggests that different knowledge types naturally
diverge at different depths, consistent with this hypothesis. Full
validation requires training multi-level forests with variable fork depths.

The tree also resembles a broader pattern seen in biological systems: shared
early processing followed by progressively more specialized downstream
computation. We do not claim a mechanistic correspondence to cortical
organization, but the analogy is suggestive: both biological and artificial
systems appear to benefit from hierarchical specialization under resource
constraints~\cite{bena2025modularity}.

% ============================================================================
\section{Method: Progressive LoRA Forking}
\label{sec:method}

The \tok{} architecture has three components: a shared trunk, a branching
mechanism, and a divergence loss. We describe each briefly; mathematical details
are in Appendix~\ref{app:math}.

\subsection{LoRA on FFN Layers, Not Attention}

A key design choice: we apply LoRA adapters to the \emph{feed-forward network}
(specifically \texttt{gate\_proj} and \texttt{up\_proj}) rather than to
attention layers. While Hu et al.~\cite{hu2022lora} demonstrated LoRA
primarily on attention projections, mechanistic interpretability research
suggests that FFN layers function as key-value memories that store factual
associations~\cite{geva2021transformer}, whereas attention layers perform
routing and compositional binding. Our goal is to modularize knowledge---whether that modularization ultimately
follows domain boundaries (medical vs.\ legal) or functional boundaries
(structure vs.\ content) is an empirical question. The FFN is the natural
target because we are modifying what the model \emph{stores}, not how it
\emph{routes} information. The attention layers remain shared across all
experts, preserving general reasoning and compositional abilities (fluid
intelligence).

\subsection{Shared Trunk, Branching Depth}

The trunk boundary is determined empirically by layer divergence analysis
(Section~\ref{sec:divergence}). For Qwen3-1.7B (28 layers), domain-specific
signals emerge at layer~18 (CKA increases 8$\times$ between layers 14 and 18).
We partition the transformer into:
\begin{itemize}
    \item \textbf{Layers 0--17 (trunk):} Shared FFN, always VRAM-resident.
          Fluid intelligence core---general reasoning, syntax, attention
          patterns. $\sim$1.2B parameters.
    \item \textbf{Layers 18--27 (expert layers):} Variable-rank LoRA adapters
          on FFN. Each adapter $\sim$45\,KB at rank-4 per layer.
\end{itemize}
Total specialist parameters: 655K for 2 experts across 10 layers, versus
800+\,MB for full expert copies in standard MoE.

\subsection{Learntropy-Driven Mitosis}

% TODO: Learntropy definition is provisional. Wozniak (SuperMemo) independently
% coined "learntropy" as the attractiveness of a signal to the learn drive —
% a NET reward signal (reward minus decoding-failure penalty) that depends on
% prior knowledge and can go negative. Our current definition (raw CE loss) is
% a simplification: it does not distinguish "productively hard" (Wozniak's
% positive learntropy) from "impossibly hard" (negative learntropy). Bimodality
% detection partially recovers this distinction, but a proper learntropy
% function should model net learning value, not raw difficulty. We are working
% toward a definition closer to Wozniak's, operationalized as a loss function.
% See: https://supermemo.guru/wiki/Learntropy
%
\textbf{Definition: Learntropy.} We use \emph{learntropy} to denote the per-token prediction difficulty experienced by a specific expert:
\begin{equation}
\mathcal{L}_{\text{learn}}(e_i, t) = -\log P(t \mid \text{context}, e_i)
\end{equation}
where $e_i$ is expert $i$ and $t$ is the predicted token. The term was independently coined by Wozniak~\cite{wozniak2018learntropy} to denote the attractiveness of a signal to the learn drive system---a net reward that depends on prior knowledge and can go negative when material exceeds the learner's capacity. Our operationalization as raw cross-entropy loss is a simplification: it measures difficulty but does not distinguish productive difficulty (material in the learner's zone of proximal development~\cite{vygotsky1978}) from noise. Bimodality detection (below) partially recovers this distinction by identifying when an expert's loss distribution contains both learnable and unlearnable populations.

We use learntropy as a \emph{diagnostic signal for the learning system itself}: not only as a training objective but as an instrument that determines when structural change is needed. This reflects Shannon's insight that information content drives the allocation of representational resources~\cite{shannon1948}---and Piaget's observation that prediction failure (accommodation) triggers the creation of new cognitive structure~\cite{piaget1952}.

When an expert's learntropy distribution is bimodal---some tokens easy, others hard---the expert must either adapt or divide. Following biological cells, we distinguish two responses:

\textbf{Gene expression (rank increase).} The expert first attempts to handle both populations by increasing its rank---activating more parameters, like a cell expressing additional genes to handle a broader range of signals. This is the cheaper response: no routing overhead, no data partitioning.

\textbf{Mitosis (expert split).} If rank increase does not resolve the bimodality, the expert undergoes \emph{mitosis}: it divides into two daughter experts, each inheriting the parent's weights plus a small random perturbation. Like biological daughter cells, they start genetically identical but \emph{differentiate} through exposure to different environments (data partitions) and explicit pressure to diverge (contrastive loss). The procedure:
\begin{enumerate}
    \item Score each expert's tokens by learntropy.
    \item Test for bimodality (Hartigan's dip test).
    \item If bimodal and rank increase insufficient: duplicate the LoRA adapter.
          Each daughter inherits parent weights plus perturbation, starting at rank-1.
    \item Daughter experts \emph{differentiate} through their data environment and contrastive loss (Section~\ref{sec:contrastive}).
\end{enumerate}
The decision to divide is triggered by an information-theoretic signal, not a schedule---analogous to how cell division is triggered by growth signals, not a timer.

\subsection{Variable Rank as Accommodation}

The rank of an adapter measures the dimensionality of its update-subspace
relative to the parent representation. We formalize rank growth as Piagetian
\emph{accommodation}: structural expansion that occurs when new residual
structure cannot be \emph{assimilated} within the existing subspace.

For an expert $e$ with current update-subspace $U_e$ (spanned by its LoRA
matrices $A$ and $B$), let $g$ be the gradient on the residual error. Decompose
$g = g_\parallel + g_\perp$ where $g_\parallel \in U_e$ and $g_\perp \perp U_e$.
The \emph{accommodation ratio} is:
\begin{equation}
\mathcal{A}(e) = \frac{\|g_\perp\|}{\|g\|}
\label{eq:accommodation}
\end{equation}

When $\mathcal{A}(e)$ is low, the error fits within the current subspace
(assimilation---no structural change needed). When $\mathcal{A}(e)$ is high,
the error requires capacity the expert does not have. The response depends on
the \emph{structure} of $g_\perp$: if it is unimodal (one coherent residual
population), rank growth is appropriate; if it is multimodal (distinct
populations requiring different corrections), splitting is appropriate.

\textbf{Practical caveat.} At low rank, $\mathcal{A}(e)$ is dominated by
geometry: a rank-$r$ subspace in $d$ dimensions captures only $r/d$ of random
directions, yielding $\mathcal{A}(e) \approx 1 - r/d$ even for uninformative
gradients. We measured $\mathcal{A}(e) \approx 0.99$ at rank~4 across all
training phases---close to the random baseline of $1 - 4/2048 = 0.998$. A
useful trigger therefore requires either higher rank (so the subspace covers a
non-trivial fraction) or a \emph{normalized} accommodation ratio that
compares against the rank-appropriate baseline.

We distinguish two regimes:
\begin{itemize}
    \item \textbf{Rank 1--4} (4--16\,KB): A lens. Residual corrections fit
          in a small subspace; the parent representation is nearly sufficient.
    \item \textbf{Rank 16--64} (64--256\,KB): A crystal. The update-subspace
          has grown large enough to encode structure the parent cannot represent.
\end{itemize}

\begin{table}[t]
\centering
\caption{Specialization depth, adapter rank, and storage tier correlate.}
\label{tab:unification}
\begin{tabular}{@{}llll@{}}
\toprule
Specialization & Adapter & Size & Storage tier \\
\midrule
Minimal (lens) & Rank-1 LoRA & 4\,KB & Always resident \\
Moderate       & Rank-16 LoRA & 64\,KB & VRAM cache \\
Substantial (crystal) & Rank-64 LoRA & 256\,KB & NVMe (instant load) \\
Full expert    & Dense copy & 12\,MB & NVMe (prefetchable) \\
\bottomrule
\end{tabular}
\end{table}

\subsection{Contrastive Divergence Loss}
\label{sec:contrastive}

Without explicit pressure, siblings drift toward the shared optimum (we observed
CosSim 0.998). The contrastive loss provides the missing force:
\begin{equation}
\mathcal{L}_{\text{contrast}} = \lambda_c \sum_{i < j} \max\bigl(0,\; \cos(\theta_{ij}) - \tau\bigr)
\label{eq:contrastive}
\end{equation}
where $\theta_{ij}$ is the angle between sibling weight vectors and $\tau$ is a
target maximum similarity. Combined with disjoint data and LoRA initialization,
this produces the differentiation that eluded all six prior arms.

\subsection{Teacher-Guided Data Selection via Zone of Proximal Development}
\label{sec:zpd}

Contrastive loss differentiates expert \emph{weights}, but does not guarantee
that experts encode different \emph{knowledge}. Two experts can be orthogonal
in weight space yet both required for every input---as we observe with the
first-level Structure/Content split (Section~\ref{sec:hotloading}).

To produce modular experts whose removal causes \emph{selective} damage, each
expert must train on fundamentally different data. We select this data using a
teacher model and Vygotsky's Zone of Proximal Development (ZPD): the set of
examples that are learnable but not yet learned.

For each text chunk $x$ and expert $e_i$, we compute:
\begin{equation}
\text{ZPD}(x, e_i) = \frac{\text{PPL}_{\text{student}}(x, e_i)}{\text{PPL}_{\text{teacher}}(x)}
\label{eq:zpd}
\end{equation}
A high ZPD score means the student struggles where the teacher succeeds---the
knowledge exists in the teacher but is absent from the student. A low score
means either both struggle (noise, not knowledge) or both succeed (nothing to
learn). Each expert receives the chunks where \emph{its own} ZPD is highest,
producing genuinely disjoint training distributions.

We tested ZPD scoring with Qwen3-30B as teacher and Qwen3-1.7B as student on
500 C4 chunks: Spearman $\rho = 0.958$ between teacher and student difficulty
rankings, with only 14.2\% of chunks in the ZPD zone. In this configuration,
the teacher adds little signal beyond what the student already knows about its
own difficulty. This does not falsify the ZPD concept itself---the result is
specific to this model pair on generic C4 data. ZPD-guided selection may prove
more effective (a)~after experts have diverged enough to have genuinely
different difficulty profiles, (b)~with a much larger teacher, or (c)~on
domain-specific data where teacher--student knowledge gaps are wider.

% ============================================================================
\section{Experiments}
\label{sec:experiments}

All experiments use Qwen3-1.7B on C4 data. We build incrementally: each
subsection adds one capability and reports the result.

\subsection{Layer Divergence Analysis: Where to Fork}
\label{sec:divergence}

Before placing adapters, we measure where domain-specific representations
emerge in the unmodified model. We pass 200 samples from each of five C4
domains (medical, code, legal, conversational, news) through Qwen3-1.7B and
compute per-layer Centered Kernel Alignment (CKA) between domain-specific
hidden state distributions.

\begin{table}[h]
\centering
\caption{Domain divergence by layer (selected). CKA measures alignment between
domain-specific hidden state distributions (higher = more domain-specific).}
\label{tab:divergence}
\begin{tabular}{@{}lrrr@{}}
\toprule
Layer range & CKA & Fisher ratio & Interpretation \\
\midrule
0--8 (early) & 0.003--0.004 & 0.004--0.005 & Domain-agnostic \\
9--14        & 0.005--0.013 & 0.006--0.009 & Minimal divergence \\
15--17       & 0.020--0.062 & 0.010--0.012 & Transition zone \\
\textbf{18--21} & \textbf{0.106--0.200} & \textbf{0.014--0.016} & \textbf{Domain-specific} \\
22--27       & 0.206--0.216 & 0.015--0.018 & Strong divergence \\
\bottomrule
\end{tabular}
\end{table}

CKA increases by an order of magnitude between layers 14 and 18 (0.013 to
0.106). This empirically determines the trunk boundary: layers below 18 carry
insufficient domain signal for expert specialization. Our initial experiments
(below) used a trunk of 14 layers;
the layer divergence analysis motivated extending the trunk to 18 layers in
subsequent experiments.

\subsection{Baseline}

\begin{table}[h]
\centering
\caption{Qwen3-1.7B baseline on C4 validation.}
\label{tab:baseline}
\begin{tabular}{@{}lr@{}}
\toprule
Configuration & Perplexity \\
\midrule
Dense baseline (no adapters) & 21.20 \\
\bottomrule
\end{tabular}
\end{table}

The dense Qwen3-1.7B achieves a perplexity of 21.20 on the C4 validation set.
This is the number every subsequent experiment must match or improve.

\subsection{Phase 1: First Split with LoRA Adapters}

We freeze the trunk (layers 0--14) and attach rank-4 LoRA adapters to layers
15--27, training on full C4 data without domain partitioning. This validates
that the LoRA+trunk architecture can recover baseline quality.

\begin{table}[h]
\centering
\caption{Phase 1: LoRA adapters on the upper layers.}
\label{tab:phase1}
\begin{tabular}{@{}lrr@{}}
\toprule
Configuration & PPL & $\Delta$ from baseline \\
\midrule
Dense baseline & 21.20 & --- \\
Trunk + LoRA (layers 15--27) & 19.59 & $-$1.61 \\
\bottomrule
\end{tabular}
\end{table}

The LoRA adapters not only recover baseline quality but \emph{improve} it,
reducing perplexity to 19.59. The trunk preserves reasoning capacity while the
adapters specialize the upper layers.

\subsection{Phase 2: Contrastive Forking}

We create two sibling LoRA adapters at the first fork point, train them on
disjoint data partitions with the contrastive loss (Eq.~\ref{eq:contrastive}),
and measure differentiation.

\begin{table}[h]
\centering
\caption{Phase 2: Sibling differentiation with contrastive loss.}
\label{tab:phase2}
\begin{tabular}{@{}lrrr@{}}
\toprule
Configuration & CosSim & PPL & Note \\
\midrule
Disjoint data only (no contrast) & 0.975 & 19.62 & Minimal differentiation \\
Disjoint data + contrastive loss & 0.840 & 19.58 & Strong differentiation \\
\bottomrule
\end{tabular}
\end{table}

The contrastive loss reduces cosine similarity from 0.975 to 0.840---a
qualitative shift from ``nearly identical'' to ``measurably different
functions''---while \emph{improving} perplexity to 19.58. Differentiation does
not trade off against quality.

\subsection{Progressive Forking}
\label{sec:forking}

With the trunk boundary informed by the layer divergence analysis
(Section~\ref{sec:divergence}), we compare forking at two boundary points.

\begin{table}[h]
\centering
\caption{Trunk boundary comparison: Phase 2 results (2 experts, rank 4,
contrastive weight 0.1). Layer-18 trunk places experts in the domain-specific
zone identified by the divergence analysis.}
\label{tab:trunk-comparison}
\begin{tabular}{@{}lrrrr@{}}
\toprule
Configuration & Expert layers & CosSim & PPL & Params \\
\midrule
Trunk-14 (layers 15--27) & 14 & 0.840 & 19.58 & 1.9M \\
Trunk-14 + Phase 3 (rank 32) & 14 & 0.278 & 19.33 & 11.2M \\
Trunk-14 + Phase 3 (rank 4) & 14 & 0.433 & 19.45 & 1.9M \\
\textbf{Trunk-18 (layers 18--27)} & \textbf{10} & \textbf{0.771} & \textbf{20.14} & \textbf{655K} \\
Trunk-18 + Phase 3 (rank 4) & 10 & 0.311 & 19.83 & 655K \\
\bottomrule
\end{tabular}
\end{table}

The trunk-18 configuration achieves deeper differentiation (CosSim 0.311 at
rank 4) than trunk-14 at the same rank (0.433), with 3$\times$ fewer
parameters. Both fork location and rank constraint affect differentiation:
trunk-14 at rank 32 reaches 0.278, while trunk-14 at rank 4 reaches only
0.433. Trunk-18 at rank 4 matches the rank-32 result (0.276 vs.\ 0.278),
suggesting that \emph{placing experts in the domain-divergent zone compensates
for 8$\times$ fewer parameters}.

However, routing analysis (Section~\ref{sec:magnitude}) shows that none of
these configurations achieve domain-specific routing. Weight-space
differentiation (low CosSim) does not imply functional specialization. Recent
work on MoE expert roles~\cite{standing2026} confirms this pattern at scale:
a ``standing committee'' of experts handles functional/structural processing
across all domains, while domain-specific routing is peripheral.

No second fork (2$\to$4 experts) has triggered in any configuration on
generic C4 data.

\subsection{Hot-Loading Validation}
\label{sec:hotloading}

Our first hot-loading test \emph{failed}: removing either first-level expert
degraded all text types equally ($\sim$1.3 PPL each). Token routing analysis
revealed Expert~0 handles rare/content words (71\% of rare tokens) while
Expert~1 handles punctuation and function words (86\% of punctuation). Every
sentence requires both structure and content---neither expert is independently
removable.

\textbf{Observation:} The first split is functional (structure vs.\ content),
not domain-specific. Both experts are always needed.

\textbf{Resolution at 8B scale:} Domain-selective ablation damage \emph{is}
observed at 8B (Section~\ref{sec:8b}), where expert removal causes
3.5--13.5$\times$ varying damage across domains ($M_{ij}$ CV 0.50--0.67).
This does not require a second fork---causal locality emerges from the
first fork at sufficient model scale, even with near-uniform routing.

\textbf{Alternative explanations partially addressed by 8B results:}
\begin{itemize}
    \item Deeper splits may also produce functional rather than domain
          differentiation---but the 8B ablation matrix shows domain-varying
          damage (3.5--13.5$\times$), indicating that domain structure does
          emerge at sufficient scale.
    \item The structure/content split at 1.7B may be an artifact of contrastive
          loss pushing toward the easiest orthogonal direction in weight space.
          The 8B results suggest this is a scale limitation, not a fundamental one.
    \item The forest metaphor may not accurately describe the learning dynamics
          at all---the branching may reflect optimization landscape geometry
          rather than knowledge organization. This remains open.
\end{itemize}

\begin{table}[h]
\centering
\caption{Ablation matrix $M_{ij}$ for trunk-18 Phase~3 final checkpoint
(CosSim 0.311). Each cell shows $\Delta$PPL when expert $j$ is removed,
evaluated on domain $i$. A modular system would show high values on the
diagonal and low values off-diagonal.}
\label{tab:ablation-matrix}
\begin{tabular}{@{}lrr@{}}
\toprule
Domain $i$ & Remove E0 & Remove E1 \\
\midrule
Medical       & +1.38 & +1.16 \\
Code          & +1.63 & +1.47 \\
Legal         & +1.42 & +1.20 \\
News          & +1.26 & +1.05 \\
Conversational & +1.22 & +1.05 \\
\midrule
\textbf{Mean} & \textbf{+1.38} & \textbf{+1.19} \\
\bottomrule
\end{tabular}
\end{table}

The ablation matrix is uniform: coefficient of variation 0.11 across domains
(threshold for meaningful selectivity: $>$0.3). Each expert accounts for a substantial fraction of the total LoRA improvement
(high substitutability: removing either costs $\sim$1.2--1.4 PPL of the
$\sim$1.4 total benefit). This confirms that
weight orthogonality (CosSim 0.311) does not produce causal locality---the
experts are near-interchangeable general-purpose adapters, not domain
specialists.

We also tested a shared + gated-routed architecture inspired by
DeepSeekMoE~\cite{deepseek2024moe}: one always-active expert per layer plus
one gated expert with a learned sigmoid gate and sparsity penalty. The gate
converged to 0.985 (near-always-on) with domain selectivity of 0.007 (needed
$>$0.20). Given the option of using additional capacity, the model uses it
uniformly rather than selectively.

\subsection{Hessian Analysis of Routing Landscape}
\label{sec:hessian}

To determine whether uniform routing is a local minimum (basin) or a saddle
point, we computed the Lanczos eigenspectrum of the LM loss Hessian restricted
to router parameters (150 iterations per layer, float32).

\begin{table}[h]
\centering
\caption{Hessian eigenspectrum at the uniform-routing solution. Nearly half of
all sampled directions point downhill: uniform routing is a saddle point.}
\label{tab:hessian}
\begin{tabular}{@{}lrrr@{}}
\toprule
Layer & Min $\lambda$ & Max $\lambda$ & Negative fraction \\
\midrule
18 & $-0.050$ & $+0.131$ & 47\% \\
22 & $-0.062$ & $+0.084$ & 48\% \\
27 & $-0.298$ & $+0.158$ & 49\% \\
\bottomrule
\end{tabular}
\end{table}

The analysis revealed two findings. First, the \textbf{root cause} of uniform
routing: our router uses hard argmax, a non-differentiable operation.
$\partial\mathcal{L}/\partial w_{\text{router}} = 0$ exactly---the router
weights are invisible to the optimizer and cannot learn from the LM loss.
Second, when argmax is replaced with a differentiable softmax surrogate, the
Hessian shows abundant negative eigenvalues (47--49\% of sampled directions),
confirming that the LM loss \emph{does} encode routing selectivity. Replacing argmax with
Gumbel-softmax~\cite{jang2017gumbel} ($\tau$: 1.0$\to$0.1) confirms this:
routing selectivity reaches 0.41 at layer~27 (vs.\ $<$0.01 with argmax),
with 82\% of tokens routed with $>$0.3 confidence---all without PPL
degradation (20.14, identical to argmax). This achieves level-2
differentiation (routing) for the first time. However, ablation testing shows
$M_{ij}$ remains uniform (CV~0.12): the routing selectivity does not translate
to causal locality. Instead, the Hessian escape directions lead to
\emph{magnitude asymmetry}---one expert becomes 1.7$\times$ more impactful
than the other across all domains, rather than each expert owning a domain.
Level-2 differentiation (selective routing) is achievable through
differentiable routing at 1.7B; level-3 (causal locality) is achieved at
8B scale (Section~\ref{sec:8b}, $M_{ij}$ CV 0.50--0.67). A
domain-conditional experiment (4 experts with gradient masking per domain,
supervised router) also failed: the general-domain expert absorbed
$>$99\% of useful learning while domain-specific experts contributed
$<$0.03 PPL each, due to severe data imbalance in C4 ($>$80\% general text).
Even forced specialization on keyword-based domain labels cannot overcome
C4's data imbalance. A \textbf{teacher-driven niche curriculum}
using token-level ZPD scores (Section~\ref{sec:zpd}) achieves strong
\emph{routing} selectivity (92\% of high-ZPD tokens route to the niche
expert), but hot-loading validation shows the expert's contribution is
negligible: removing it causes zero PPL change. The router learned
\emph{where} to send niche tokens, but the rank-4 expert did not develop
enough capacity to become indispensable. This confirms that routing
selectivity (level~2) and causal locality (level~3) are genuinely
distinct: selective routing is achievable through curriculum design, but
selective indispensability requires sufficient expert capacity or longer
training. A top-2 routing variant
with 4 experts (Gumbel-softmax, each token activates 2 of 4) developed a
\emph{capacity hierarchy}---two ``heavy'' experts contributing $\sim$0.56 PPL
each and two ``light'' experts contributing $\sim$0.28---but no domain
selectivity (per-expert CV 0.11--0.16). This is the Standing Committee
pattern~\cite{standing2026} at 4-expert scale: functional importance
differentiates, domain ownership does not. Finally, a damage-prediction
surrogate that directly optimizes for selective indispensability
($L_\text{loc} \propto \text{overlap} - \lambda\,\text{exclusivity}$
of per-token ablation damage) achieves high token-level exclusivity
(1.27 at training time) but the resulting $M_{ij}$ is again
domain-uniform (CV~0.12): expert~1 carries 2.5$\times$ more load than
expert~0, uniformly across all domains. The experts own different
\emph{tokens}, but those tokens are distributed evenly across domains.

\subsection{LoRA Magnitude Analysis}
\label{sec:magnitude}

To understand \emph{why} hot-loading fails, we measured the magnitude of LoRA
deltas relative to the trunk FFN output across 1000 evaluation tokens.

\begin{table}[h]
\centering
\caption{LoRA magnitude analysis at layers 14--27 (Phase 3 checkpoint,
rank-32). Pre-registered prediction: ratio $<$0.05 (small correction).
Result: ratio $\approx$0.25 (substantial correction).}
\label{tab:magnitude}
\begin{tabular}{@{}lrrr@{}}
\toprule
Metric & Expected & Measured & Verdict \\
\midrule
$\|\text{LoRA}\| / \|\text{FFN}\|$ (mean) & $<$0.05 & 0.249 & FAIL \\
Gate CosSim between experts & --- & 0.019 & Orthogonal \\
Routing selectivity (pref.\ vs non-pref.) & $>$2$\times$ & $<$5\% diff & FAIL \\
\bottomrule
\end{tabular}
\end{table}

The LoRA adapters contribute $\sim$25\% of the FFN output magnitude---they
are not small corrections but substantial transformations. Expert weight
matrices are nearly orthogonal (CosSim 0.019), confirming that contrastive
loss produces genuine weight-space differentiation. However, routing
selectivity is weak: both experts produce similar-magnitude deltas regardless
of which tokens are routed to them. This reveals \textbf{redundant representation}: the experts learned different
directions in representation space, but both directions are needed by most
tokens.

We initially hypothesized that this was caused by placing experts in a
domain-agnostic zone (layers 14--17), motivated by the layer divergence
analysis (Section~\ref{sec:divergence}). However, repeating the routing
analysis on the trunk-18 configuration (experts on layers 18--27, where CKA
$>$0.10) yielded the same result: both experts route 49--51\% of tokens
uniformly across all five domains, with per-expert PPL ablation varying by
only $\sim$0.26 points across domains. \textbf{Moving the fork to the
domain-divergent zone did not produce domain-specific routing.}

This suggests that the redundancy is not caused by fork location but by a more
fundamental limitation: contrastive loss differentiates expert \emph{weights}
but does not create functional \emph{routing specialization}. The router
learns to split tokens approximately 50/50 regardless of domain content,
and both experts contribute equally to all domains. Weight-space orthogonality
(CosSim $<$0.3) is a necessary but not sufficient condition for modularity.

% ============================================================================
\section{Related Work}
\label{sec:related}

\paragraph{LoRA and adapters.}
Hu et al.~\cite{hu2022lora} introduced Low-Rank Adaptation for efficient
fine-tuning. We extend LoRA from a fine-tuning tool to a \emph{structural unit
of specialization} within a single model. Li et al.~\cite{li2018intrinsic}
and Aghajanyan et al.~\cite{aghajanyan2021intrinsic} established that
fine-tuning operates on a low intrinsic dimension---our variable-rank mechanism
exploits this directly. HydraLoRA~\cite{hydralora2024} demonstrates that asymmetric residual
specialization (shared A + routed B matrices) is viable, though it does not
test causal modularity (hot-pluggability). MoLoRA~\cite{molora2026} validates
per-token LoRA routing and composability independently.

\paragraph{Shared vs.\ routed experts.}
DeepSeekMoE~\cite{deepseek2024moe} introduced \emph{shared expert isolation}:
always-active experts absorb cross-domain knowledge, freeing routed experts to
specialize. DeepSeek-V3~\cite{deepseekv3} scales this to 1 shared + 256
routed experts. The ``Basic-Refinement Collaboration''
framework~\cite{basicrefinement2025} empirically confirms that shared experts
perform entity recognition and syntactic parsing while routed experts refine
domain-specific associations. Our finding that experts differentiate along
functional axes but not domain axes (Section~\ref{sec:magnitude}) is
independently reported by OpenMoE~\cite{openmoe2024} (context-independent
token-type routing), the Mixtral analysis~\cite{mixtral2024} (``no obvious
patterns based on topic''), and the Standing
Committee~\cite{standing2026} paper (functional experts as stable core, domain
experts as peripheral). The ``Domain vs.\ Driver Experts''
analysis~\cite{domaindriver2026} formalizes this: driver experts handle
function words and syntax, domain experts handle content vocabulary.

\paragraph{Specialization theory.}
B\'{e}na \& Goodman~\cite{bena2025modularity} prove resource constraints are
necessary for specialization---without cost pressure, modules converge. Our
contrastive loss provides one form of constraint; the rank measures the result.
However, our experiments show that weight-space constraints (contrastive loss)
produce representational but not functional specialization. The literature
suggests that top-$K$ routing ($K > 1$) may be required: with $K=1$, each
token must choose a single axis of specialization (typically functional), while
$K \geq 2$ allows simultaneous functional and domain expert
activation~\cite{standing2026, domaindriver2026}.

\paragraph{Low-rank routing.}
L2R~\cite{l2r2026} projects to $r{=}2$ dimensions for routing, confirming that
routing decisions are intrinsically low-dimensional. Our variable-rank adapters
extend this: the routing is low-rank, and so is the specialization itself.

% ============================================================================
\section{Analysis}
\label{sec:analysis}

We address two questions about the emergent structure, reformulated from
our original framing in light of the 2-expert results.

\paragraph{Q1: How does rank affect information efficiency?}
We measure the residual compression ratio (RCR = PPL gain / trainable
parameters) across ranks 1--32 at the trunk-18 boundary.
\begin{table}[h]
\centering
\caption{Rank sweep at trunk-18 boundary (5K steps each, layers 18--27).
Gain measured against frozen-trunk baseline (PPL 26.1, no adapters).
RCR = gain / trainable parameters in millions.}
\label{tab:rank-sweep}
\begin{tabular}{@{}rrrrr@{}}
\toprule
Rank & Params (M) & PPL & Gain & RCR \\
\midrule
1  & 0.164 & 20.73 & 5.42 & \textbf{33.1} \\
2  & 0.328 & 20.65 & 5.49 & 16.8 \\
4  & 0.655 & 20.61 & 5.54 & 8.4 \\
8  & 1.311 & 20.55 & 5.60 & 4.3 \\
16 & 2.621 & 20.43 & 5.71 & 2.2 \\
32 & 5.243 & 20.29 & 5.85 & 1.1 \\
\bottomrule
\end{tabular}
\end{table}

RCR drops 30$\times$ from rank-1 to rank-32, showing that the first few
dimensions of the update-subspace capture most of the learnable signal. In
absolute terms, even rank-32 adds only 5.2M parameters (0.26\% of the 2B base
model) for a 0.44 PPL improvement over rank-1---still highly efficient
relative to the base model. The rank sweep thus validates both the lens and
crystal regimes: low-rank adapters are maximally information-efficient per
parameter, while high-rank adapters remain negligible in total model storage
and can be justified when the marginal quality matters.

\paragraph{Q2: Is the emergent specialization interpretable?}
The first fork produces a structure/content split: Expert~0 routes 71\% of
rare words, Expert~1 routes 86\% of punctuation. This is functional, not
domain-based. The split aligns with the ``Standing Committee''
pattern~\cite{standing2026}: functional experts form a stable core, domain
experts are peripheral. At Qwen3-1.7B scale, domain-level interpretability
has not emerged---but at 8B, ablation damage varies 3.5--13.5$\times$
across domains (Section~\ref{sec:8b}), indicating that domain-level causal
structure does emerge at sufficient scale.

% ============================================================================
\section{Limitations}
\label{sec:limitations}

We identify four limitations that bound our current claims:

\begin{enumerate}
    \item \textbf{Weight differentiation does not imply functional
          specialization at 1.7B.} This is the central negative finding at
          small scale, resolved by the 8B experiments.
          Contrastive loss reliably produces orthogonal expert weights
          (CosSim $<$0.3) at both fork boundaries tested (layers 14 and 18).
          However, token routing analysis shows both experts route all domains
          uniformly (49--51\% split), and removing either expert degrades all
          domains equally. The experts learn different \emph{directions} in
          representation space, but both directions are needed by all tokens.
          At 1.7B, causal modularity is not achieved ($M_{ij}$ CV $\leq$ 0.16).
          However, scaling to 8B resolves this: $M_{ij}$ CV reaches 0.50--0.67
          with domain-selective ablation damage (Section~\ref{sec:8b}).
          Hessian analysis (Section~\ref{sec:hessian}) shows the LM loss
          \emph{does} encode selective routing directions; sufficient model
          capacity is required to exploit them.
    \item \textbf{Split timing is scheduled, not learntropy-driven.}
          The paper's theoretical framework positions learntropy as the
          signal for \emph{when} to split. In practice, all successful
          multi-expert experiments use a fixed schedule (split every $N$
          steps). Three candidate signals were tested retroactively:
          intra-expert loss variance (ILV), router margin bimodality,
          and learntropy trend (slope). All three correctly predict
          \emph{which} expert should split (the one with highest
          learntropy), but none provides a sharp \emph{when} signal---all
          experts' learntropy trends downward simultaneously, with no
          clear plateau or divergence preceding scheduled split points.
          Bimodality detection (dip test, kurtosis) on per-token loss
          distributions never fires on generic C4 data. Finding a
          mechanistic, data-driven split trigger remains an open problem.
          The scheduled approach works well in practice. A promising
          alternative is \emph{speculative splitting with rollback}:
          attempt a split at each candidate point, train for a short
          evaluation window, and revert if the split did not reduce loss.
          This inverts the problem---instead of predicting \emph{when}
          to split, the system tries splitting and evaluates whether it
          helped. In a 25K-step experiment, the system accepted only
          2 out of 8 attempted splits (at steps 9K and 21K), growing
          to 4 experts with all active ($M_{ij}$ CV 1.52). Early splits
          (3K, 6K) were correctly rejected as premature. The cost is at
          most the evaluation window (500 steps) per failed attempt---a
          modest overhead that eliminates the need for any timing signal.
    \item \textbf{Framework tested on one model family only.} All experiments use
          Qwen3 (1.7B and 8B). Generalization to other architectures (Llama,
          Mistral, GPT-scale) and other data distributions is assumed but not
          demonstrated.
    \item \textbf{Learntropy definition is provisional.} Our use of raw
          cross-entropy as a splitting signal does not distinguish productively
          difficult tokens from noise (Section~\ref{sec:method}). A principled
          learntropy function that models net learning value---closer to
          Wozniak's~\cite{wozniak2018learntropy} definition---may be needed
          for effective curriculum-driven forking.
\end{enumerate}

% ============================================================================
\section{Scale Experiment: Qwen3-8B}
\label{sec:8b}

All 1.7B experiments share a common ceiling: $M_{ij}$ CV $\leq 0.16$ regardless
of configuration. To test whether this is a scale limitation, we replicate the
LoRA forking experiment on Qwen3-8B (36 layers, 4096 hidden, 11008
intermediate). Layer divergence analysis identifies layer~24 as the fork
boundary (fisher ratio 0.081 vs.\ 0.018 at 1.7B), and we train 2-expert
rank-4 LoRA adapters on layers 24--35 for 25K steps on C4.

\paragraph{Causal locality achieved.}
Table~\ref{tab:8b-mij} shows the ablation matrix. Expert~0 removal causes
2--5$\times$ more damage than Expert~1, and damage varies 3.5--13.5$\times$
across domains. Conversational text is most vulnerable; news is least. The
$M_{ij}$ CV reaches 0.50--0.67 across three independent runs (5K, 25K steps,
and a 4-expert variant)---\textbf{the first domain-selective causal locality
observed in our experiments}.

\begin{table}[h]
\centering
\caption{8B ablation matrix ($\Delta$PPL when expert removed). Domains show
3.5--13.5$\times$ varying damage---causal locality.}
\label{tab:8b-mij}
\begin{tabular}{@{}lrrrrr@{}}
\toprule
 & Medical & Code & Legal & News & Conv. \\
\midrule
$-$Expert 0 & +1.69 & +2.00 & +2.50 & +1.25 & +4.63 \\
$-$Expert 1 & +0.69 & +0.88 & +1.13 & +0.31 & +2.13 \\
\midrule
Ratio (E0/E1) & 2.5$\times$ & 2.3$\times$ & 2.2$\times$ & 4.0$\times$ & 2.2$\times$ \\
\bottomrule
\end{tabular}
\end{table}

\paragraph{Quality benchmark.}
We compare the 8B expert-augmented model against the unmodified base on
perplexity across 30 C4 validation samples per domain
(Table~\ref{tab:8b-quality}). The LoRA experts improve perplexity by
16--22\% across \emph{all} domains, confirming that the adapters add genuine
knowledge without degrading the base model's capabilities.

\begin{table}[h]
\centering
\caption{PPL comparison: base Qwen3-8B vs.\ FoK-8B (with 2 LoRA expert
adapters, rank-4, layers 24--35). 30 samples per domain from C4 validation.}
\label{tab:8b-quality}
\begin{tabular}{@{}lrrr@{}}
\toprule
Domain & Base PPL & FoK PPL & $\Delta$ \\
\midrule
Medical       & 14.35 & 12.06 & $-$16.0\% \\
Code          & 23.56 & 18.93 & $-$19.6\% \\
Legal         & 19.13 & 14.98 & $-$21.7\% \\
News          & 14.20 & 11.66 & $-$17.9\% \\
Conversational& 23.75 & 18.98 & $-$20.1\% \\
\midrule
\textbf{Overall} & \textbf{18.52} & \textbf{14.99} & $-$\textbf{19.1\%} \\
\bottomrule
\end{tabular}
\end{table}

\paragraph{Routing vs.\ causality.}
Routing weights at inference are nearly uniform (47/53 E0/E1 split across all
domains), yet ablation damage varies strongly. This confirms the three-level
framework: level-2 (routing selectivity) and level-3 (causal locality) are
independent. The experts are causally specialized in \emph{what they compute},
not in \emph{which tokens they process}---an important distinction from
standard MoE where routing determines functionality.

\paragraph{4-expert variant: natural collapse to 2 effective experts.}
A 4-expert rank-4 configuration trained simultaneously reveals that only 2 of 4
experts learn meaningful weights. Expert~0 removal is catastrophic
($\Delta$PPL +87 to +169 across domains), Expert~1 removal is significant
(+10 to +17), while Experts~2 and~3 contribute zero ($\Delta$PPL $<$0.02).
Crucially, routing weights are perfectly uniform (25\% per expert) despite
a 10$\times$ asymmetry in causal importance. Progressive forking (training 2
experts to convergence, freezing them, then adding 2 new experts) produces
the same collapse: the frozen experts dominate so thoroughly that new experts
receive no useful gradient. This provides the strongest evidence that routing
weights are \emph{uninformative} about functional specialization---the
three-level separation is not an artifact of 2-expert configurations but a
fundamental property of how LoRA experts differentiate.

The collapse to 2 effective experts despite 4 available slots suggests a
deeper problem: all experts share the same global learning rate, so early
learners monopolize gradient signal. In Piaget's
framework~\cite{piaget1952}, this is the distinction between
\emph{assimilation} (adjusting input to existing schemas) and
\emph{accommodation} (restructuring schemas for new input). A discrete
freeze/train cycle is a crude approximation of what should be a continuous
process.

\paragraph{Early fork (layer 12): testing variable override depth.}
We replicate the 2-expert experiment with LoRA adapters on layers 12--35
instead of 24--35, doubling the number of expert layers (24 vs.\ 12) and
trainable parameters. The results show extreme asymmetry: Expert~0 removal
causes catastrophic damage ($\Delta$PPL +7,676 to +15,955), while Expert~1
removal is mild (+5.6 to +7.7). Expert~0 absorbed essentially all learning
capacity.

Critically, the domain damage pattern is consistent with the forest
hypothesis: \textbf{code receives the most damage} when Expert~0 is removed
(+14,461) and \textbf{medical the least} (+7,676)---a 1.88$\times$ ratio.
This matches the prediction that code diverges earlier from natural language
than medical text. The overall $M_{ij}$ CV is 1.07 (vs.\ 0.68 at layer~24),
confirming that earlier overrides produce more domain-selective experts.

\paragraph{The rank spectrum: from lens to full layer.}
The distinction between ``LoRA adapter'' and ``full layer replacement'' is
not categorical but continuous. An FFN layer in Qwen3-8B has dimensions
$4096 \times 11008$; the maximum LoRA rank is $\min(4096, 11008) = 4096$. At
rank~4096, $A \times B$ is a full matrix---mathematically equivalent to a
complete layer replacement. The accommodation ratio $A(e)$ determines where
each expert sits on this spectrum:

\begin{center}
\begin{tabular}{@{}rrl@{}}
\toprule
Rank & \% of full layer & Character \\
\midrule
4 & 0.1\% & Lens (redirects information) \\
64 & 1.6\% & Crystal (encodes own features) \\
256 & 6.3\% & Semi-autonomous \\
4096 & 100\% & Full layer replacement \\
\bottomrule
\end{tabular}
\end{center}

This unifies the storage-tiering story: rank-4 adapters (4\,KB) are always
cached, rank-4096 ``adapters'' ($\sim$32\,MB) load from NVMe. Same API, same
routing, same mechanism---only the matrix dimensions change.

\paragraph{Rank-8 confirms the information-theoretic ceiling.}
Doubling the rank from 4 to 8 at layer~24 (with the same learntropy-LR
mechanism) increases the $M_{ij}$ CV from 0.68 to 1.09 (2 experts, +61\%).
With 4 experts at rank-8 and learntropy-LR, 3 of 4 experts are active and
the overall $M_{ij}$ CV reaches \textbf{1.72}---the highest observed in any
experiment. The assimilation expert (lowest learntropy, lowest LR) again
becomes the most critical ($\Delta$PPL +249 to +421), while two
accommodation experts go dormant. This confirms that the 2-expert
equilibrium at rank-4 is a consequence of operating in too few dimensions,
and that rank and learntropy-LR are complementary: rank provides the
subspace, learntropy-LR distributes the learning signal.

However, this relationship is \textbf{non-monotonic}. A rank sweep at
layer~12 reveals a peaked curve:

\begin{center}
\begin{tabular}{@{}rrrl@{}}
\toprule
Rank & \% hidden & $M_{ij}$ CV & Regime \\
\midrule
4 & 0.1\% & 1.07 & LoRA (plateau) \\
16 & 0.4\% & 1.07 & MeRA (plateau) \\
64 & 1.6\% & \textbf{1.19} & Crystal (peak) \\
256 & 6.3\% & 0.77 & Semi-autonomous (declining) \\
1024 & 25\% & 0.47 & HiRA (overcapacity) \\
\bottomrule
\end{tabular}
\end{center}

Domain selectivity \emph{peaks} at rank-64 (1.6\% of hidden dimension),
then drops sharply. At rank-64, experts have just enough capacity to encode
domain-specific features but not enough to learn everything---the
information bottleneck is tight enough to force specialization but loose
enough to permit it. At rank-1024, the bottleneck disappears: each expert
can represent all domains, and differentiation pressure vanishes. A
variable-rank configuration (rank-4 on early layers, rank-16 on middle,
rank-64 on late) produces \emph{lower} selectivity (CV$=$0.62) than any
uniform rank---the early bottleneck restricts information flow while the
late high-rank layers compensate, netting less differentiation than
uniform rank throughout.

\paragraph{Hierarchical bifurcation: 4/4 experts active.}
All previous 4-expert experiments used flat 1$\to$4 routing: every expert
branches from the same layer. The original tree design uses progressive
bifurcation: 1$\to$2 at layer~18 (pair-level, functional) then 2$\to$4 at
layer~24 (leaf-level, domain). Experts within each pair share LoRA adapters
on layers 18--23 before diverging at layer 24.

This produces \textbf{4 of 4 active experts}---the first time all experts
contribute---with $M_{ij}$ CV of 1.62. The hierarchy creates a natural
gradient of importance: Expert~3 (Pair~B, leaf) is the backbone
($\Delta$PPL +116 to +178), Expert~2 (Pair~B, leaf) is significant (+5.7
to +10.2), and Experts~0/1 (Pair~A) provide modest but genuine
contributions (+0.6 to +1.1). Pair-level ablation confirms the asymmetry:
removing Pair~B causes 3.6$\times$ more damage than removing Pair~A.

The hierarchical structure succeeds where flat routing fails because
\emph{within-pair weight sharing forces differentiation at both levels}.
Pair-level experts must differentiate from each other (functional split),
and leaf-level experts within a pair must further differentiate (domain
split). This two-stage constraint prevents the collapse to $<$4 effective
experts that occurs with flat 1$\to$4 routing.

A deeper 1$\to$2$\to$4$\to$8 tree (layers 12/18/24) collapses to only 2 of
8 active leaf experts---both in the same sub-pair. The three-level hierarchy
dilutes gradient through three routing stages, and only the ``winning path''
accumulates meaningful weights. Two bifurcation levels appears to be the
sweet spot for 15K training steps at this scale.

\paragraph{Hierarchy and rank are both necessary.}
A rank-4 version of the 2-level tree produces only 2 of 4 active
experts---the same count as flat rank-4 routing. Pair~A experts (rank-4,
layers 18--35) go dormant: the 4-dimensional subspace provides too few
directions for experts within a pair to differentiate. At rank-8, the same
hierarchy activates all 4 experts. This establishes that \textbf{both
hierarchical structure and sufficient rank are required}: the hierarchy
provides the differentiation constraint (pair-level vs.\ leaf-level),
while the rank provides the representational capacity for that constraint
to produce distinct experts.

\paragraph{Scheduled growth: 8/8 experts from learntropy-directed splitting.}
A simpler growth strategy succeeds where bimodality detection fails:
every 3,000 steps, split the expert with the highest learntropy (the one
struggling most). Starting from 2 experts at rank-8 (layer~12), the tree
grows to 8 experts over 20K steps. \textbf{All 8 experts are active},
with $M_{ij}$ CV of 1.98---the highest observed in any experiment.

The growth produces a natural genealogy: Expert~0 (lowest learntropy,
never split) becomes the assimilation backbone ($\Delta$PPL +35 to +110).
Expert~1 (highest early learntropy) splits three times, producing
Experts~2, 3, and~4. Expert~4 inherits high learntropy and becomes a
secondary branch point, producing Experts~5, 6, and~7. The tree grows
asymmetrically---branches emerge where the model struggles, not on a
predetermined schedule. A reproducibility run (seed=42) confirms the
pattern: 8/8 active, $M_{ij}$ CV of 2.34 (vs.\ 1.98 in run~1). The growth
dynamics are robust: Expert~0 never splits in either run, the
accommodation lineage branches predictably, and all 8 experts contribute
to the final model. When combined with learntropy-driven rank increases
(starting at rank-4, doubling when learntropy plateaus), the system
autonomously grows from 2~experts at rank-4 to 8~experts at rank-64
through 4~rank doublings and 6~splits---converging to the empirically
optimal rank from the rank-selectivity sweep without explicit guidance.
A quality benchmark confirms that the 8-expert model reduces perplexity by
22.8\% across all domains compared to the base---an improvement over the
2-expert model's 19.1\%, with legal text benefiting most ($-$25.4\%).

\paragraph{Grove of Trees: combining tree growth with tree-level identity.}
The scheduled growth mechanism operates within a single routing space. A
\emph{grove} extends this by giving each tree its own identity: LoRA
adapters on layers~0--11 define the tree's ``personality,'' while
layers~12--35 host internal experts that grow via scheduled splits. Two
trees, each growing to 4~internal experts (8~total), produce 8/8 active
leaf experts with $M_{ij}$ CV of 1.34 and inter-tree cosine similarity of
0.044 (near-orthogonal). Tree~0 serves as the dominant computation path
($\Delta$PPL +2.5 to +43 when internal experts removed), while Tree~1
provides orthogonal corrections ($\Delta$PPL +0.1 to +3.0). The
contrastive loss between trees was ineffective (trees naturally diverge
to CosSim $<$0.05 without explicit pressure), suggesting that the
grove architecture itself enforces sufficient differentiation.

\paragraph{Zipf-distributed ablation damage enables natural tiered storage.}
The spectrum of per-expert ablation damage is strongly Zipf-distributed: a
small number of root-level adapters deliver large gains ($\Delta$PPL up to
+43 for the primary expert in Tree~0), while the long tail of leaf experts
contributes only +0.1 to +3.0 each. This skew is not a limitation but the
key enabler of efficient tiered deployment.

Expert relevance is inherently input-dependent and is determined on-the-fly
by the Gumbel-softmax router. Consequently, no separate cache policy is
required. High-damage adapters (frequently activated, encoding common
patterns) remain permanently resident in VRAM. Mid-tier adapters are kept in
the fast NVMe/flash cache, while the long tail of low-damage, low-rank
specialists (4--256\,KB) is loaded on demand from flash storage only when
the router assigns them non-negligible probability. A lightweight prefetcher
that monitors router probabilities is sufficient to keep latency negligible.

Because adapter rank---governed directly by the accommodation
ratio---already encodes the appropriate storage tier, the architecture
aligns naturally: rank-4 adapters (4\,KB) incur near-zero cost even if
kept resident, while rank-64 adapters (256\,KB) are loaded only for their
specific domain or topic. The resulting Zipf-like distribution of adapter
sizes matches the Zipf-like access patterns of real-world knowledge
(Appendix~\ref{app:shannon}).

\paragraph{Interpretability: functional, not domain-selective routing.}
Routing analysis of the 4/4-active rank-8 tree reveals that pair-level
routing is 97\% Pair~B, 3\% Pair~A---\emph{identical across all domains}.
The domain routing spread is 0.002, confirming that the pair split is
\textbf{functional} (structure vs.\ content), not domain-based. Token-type
analysis shows punctuation tokens route 0\% to Pair~A while alphabetic
tokens route 3\%, consistent with a content$\to$A / structure$\to$B
pattern. Pair-level weight cosine similarity is $-0.10$ (orthogonal).

This definitively validates the three-level framework: routing weights
(level~2) are domain-blind, but ablation damage (level~3) varies
3.5--13.5$\times$ across domains. The experts are causally specialized
through \emph{what they compute}, not \emph{which tokens they receive}.
The tree hierarchy produces differentiation through shared-weight
constraints, not through selective routing.

% ============================================================================
\section{Domain-Selective Gating: Single-Adapter Deployment}
\label{sec:delta-gating}

The multi-expert experiments in Section~\ref{sec:8b} use C4 generic data with
simultaneously-trained experts. A separate question is whether a \emph{single}
domain-specific adapter can be gated to activate selectively---high on
in-domain text, low on out-of-domain. This is the deployment mechanism for
hot-pluggable adapters: a gate that turns on when the adapter is relevant and
off when it is not.

\paragraph{Architecture.}
A per-layer scalar gate $g_l$ controls the adapter delta:
\begin{equation}
    y = f_{\text{base}}(x) + \sigma(g_l(x)) \cdot \bigl(f_{\text{adapter}}(x) - f_{\text{base}}(x)\bigr)
\end{equation}
where $g_l$ is a learned linear projection with bias initialized to $-2.0$
($\sigma(-2) \approx 0.12$, starting near OFF). Training is two-phase: (1)~train
the adapter on domain data with gate disabled, then (2)~freeze the adapter and
train only the gates on a 50/50 mix of domain and generic data. A sparsity
penalty on gate activations pushes gates toward zero by default.

\paragraph{Selectivity is robust (4 seeds).}
Across 4 random seeds (42, 256, 1337, 7), the gate consistently differentiates
domain from generic text: mean selectivity (domain gate $-$ generic gate)
$+0.621 \pm 0.009$. Domain gates converge to 0.94--0.96, generic gates to
0.28--0.34. The variance across seeds is remarkably low.

\paragraph{Generic text improves.}
Counterintuitively, the gated adapter \emph{improves} generic (non-domain)
perplexity by $-6.0\% \pm 0.1\text{pp}$ relative to the base model (4 seeds,
all improving). The gate acts as a soft regularizer: by selectively applying
the adapter delta, it provides a small benefit even on text the adapter was not
trained on.

\paragraph{Per-layer gate structure is meaningful.}
Zeroing individual layer gates produces non-uniform PPL impact. The correlation
between gate magnitude and domain PPL impact is strong (Spearman $\rho = 0.717$,
$p = 8.0 \times 10^{-5}$), while the correlation with generic PPL impact is
weak ($\rho = 0.21$, $p = 0.33$), confirming that gates are domain-specific.
The highest-gate layer (L35, gate 0.993) causes $+4.91\%$ domain PPL when
ablated; the lowest (L30, gate 0.127) causes only $+0.35\%$.

\paragraph{Cross-domain generalization.}
The same architecture and hyperparameters produce domain-selective gating on a
second, unrelated domain (Dutch cuisine): selectivity $+0.420$, domain PPL
$-64.2\%$. A 2$\times$2 cross-evaluation matrix shows diagonal dominance:
each adapter strongly improves its own domain (BBC $-73.5\%$, cuisine $-64.2\%$)
and modestly or not the other (BBC adapter on cuisine $-4.5\%$, cuisine adapter
on BBC $-13.7\%$).

\paragraph{Sparsity mechanism is not load-bearing.}
L1 and L2 penalties produce identical selectivity (gap $+0.002$) and PPL
($-17.2\%$ vs.\ $-17.4\%$). Dropout on gate activations produces lower
selectivity ($+0.525$) but better domain PPL ($-20.3\%$). The choice of
sparsity mechanism does not materially affect the outcome; any pressure toward
gate sparsity is sufficient.

\paragraph{Gate initialization is not fragile.}
The gate bias init controls the starting point on the sigmoid curve
($\sigma(-1) = 0.27$ to $\sigma(-4) = 0.018$). Across a 10$\times$ range
of initial gate activation ($-1.0$ to $-4.0$), selectivity varies only 0.025
($+0.607$ to $+0.632$) and PPL varies $<$0.2\%. The mechanism is robust to
initialization.

\paragraph{Limitation: independently trained gates do not fully compose.}
When two adapters (BBC news + Dutch cuisine) are loaded simultaneously with
independent gates, each gate is highest on its own domain. However, the
cuisine gate fires at 0.49 on BBC text (above the 0.40 ideal threshold),
adding a small interference cost ($+5.9\%$ and $+6.8\%$ PPL vs.\ single
adapter). Independently trained gates partially compose but may require
joint training or cross-adapter normalization for clean multi-adapter
deployment.

% ============================================================================
\section{Conclusion}

We set out to separate fluid intelligence from crystallized knowledge using
variable-rank LoRA expert adapters. The results establish five principal
findings and a unifying framework driven by learntropy.

\subsection{Established Results}

\textbf{1.~Information efficiency.} Variable-rank LoRA adapters on FFN layers
reduce perplexity from 21.20 to 19.83 with 655K trainable parameters on
1.7B, and by 19\% on 8B. Adapters of $\sim$45\,KB per layer load from
NVMe in microseconds. The rank is a measurement: low rank = assimilation
(lens-like correction), high rank = accommodation (new capacity).

\textbf{2.~Causal locality as the correct metric.} Weight orthogonality,
routing selectivity, and causal locality (selective ablation damage) are three
genuinely distinct levels. All 1.7B configurations achieve levels 1--2 but not
3 ($M_{ij}$ CV $\leq$ 0.16). Scaling to 8B breaks through: $M_{ij}$ CV
0.50--0.67, with ablation damage varying 3.5--13.5$\times$ across domains.
Routing weights remain nearly uniform---causal locality arises from
\emph{what} experts compute, not \emph{which tokens} they process.

\textbf{3.~The 4-expert collapse.} Training 4 experts simultaneously or
progressively produces the same result: 2 active experts, 2 dead. Routing is
perfectly uniform (25\% each) despite 10$\times$ causal asymmetry. Rank-4
adapters in a 4096-d space correct only 0.1\% of available directions---the
residual has $\sim$2 dominant components, an information-theoretic ceiling.

\textbf{4.~Learntropy-directed growth produces 8 active experts.}
Splitting the highest-learntropy expert every $N$ steps grows a tree from 2
to 8 experts with all 8 contributing ($M_{ij}$ CV up to 2.34, reproduced
across seeds). The assimilation backbone (lowest learntropy) never splits;
the accommodation lineage branches repeatedly. The 8-expert model reduces
perplexity by 22.8\% compared to the base---an improvement over the
2-expert model's 19.1\%.

\textbf{5.~Autonomous rank discovery.} When learntropy-driven rank increases
are added (doubling rank when learntropy plateaus), the system autonomously
grows from rank-4 to rank-64 through 4~doublings---converging to the
empirical peak of the rank-selectivity curve without explicit guidance.
Domain selectivity peaks at rank-64 (1.6\% of hidden dimension) and
\emph{declines} at higher ranks as overcapacity eliminates specialization
pressure.

\textbf{6.~Domain-selective gating for deployment.} A per-layer scalar gate
on the adapter delta activates selectively: domain gate 0.94--0.96, generic
gate 0.28--0.34, selectivity $+0.621 \pm 0.009$ (4 seeds). The gated adapter
improves generic perplexity by $-6.0\% \pm 0.1\text{pp}$ relative to the base
model. Per-layer gate magnitudes correlate strongly with domain PPL impact
(Spearman $\rho = 0.717$, $p < 10^{-4}$) but not with generic PPL
($\rho = 0.21$, $p = 0.33$). The mechanism generalizes across domains (2
domains tested, same hyperparameters) and is robust to the choice of sparsity
penalty (L1 $\approx$ L2, gap $+0.002$).
(Section~\ref{sec:delta-gating})

\subsection{From Tree to Forest}

The layer divergence profile across 8B (Section~\ref{sec:8b}) reveals a
U-shaped pattern: high fisher ratios at layers 0--4 and 24+, low in the
middle (8--20). This suggests that different knowledge types diverge at
different depths:

\begin{itemize}[nosep,leftmargin=1em]
\item \textbf{Code $\to$ shrub}: forking early (layer 4--8), because the
      computational substrate---syntax, types, control flow---differs
      fundamentally from natural language. Broad crown, short trunk.
\item \textbf{Medical/legal $\to$ tree}: forking late (layer 24+), because
      the reasoning pattern is shared (symptom$\to$diagnosis $\approx$
      fact$\to$conclusion). Long shared trunk, narrow crown with
      domain-specific vocabulary.
\item \textbf{Conversational $\to$ grass}: little or no crystallized
      knowledge; pure fluid intelligence needs no expert fork.
\end{itemize}

This is not a single tree of knowledge but a \emph{forest}: each knowledge
type grows its own branching structure, with variable fork depth determined
by when its representations diverge from the shared trunk.

\subsection{Learntropy as Unified Control Signal}

The learntropy signal---per-token cross-entropy weighted by routing
contribution---drives all adaptation decisions:

\begin{enumerate}[nosep,leftmargin=1em]
\item \textbf{When to fork}: bimodal learntropy = the expert's schema no
      longer fits a single distribution.
\item \textbf{Where to fork}: per-layer learntropy divergence identifies the
      optimal fork boundary for each knowledge type.
\item \textbf{How fast to learn}: per-expert learntropy modulates learning
      rate---low learntropy $\to$ assimilation (LR $\downarrow$), high
      learntropy $\to$ accommodation (LR $\uparrow$).
\item \textbf{When to abstain (hypothesis)}: at inference time, the gap
      between trunk-learntropy (without adapter) and adapter-learntropy
      (with adapter) may indicate whether the model possesses relevant
      knowledge. If validated, this would reframe hallucination from an
      alignment problem to a routing problem. This remains untested:
      learntropy may correlate with token difficulty rather than factual
      incorrectness, and a reliable threshold requires empirical
      calibration.
\end{enumerate}

A learntropy-modulated per-expert LR experiment on 8B with 4 experts
produces a striking result: \textbf{3 of 4 experts become active}---the
first time more than 2 experts have survived training in our framework
(Table~\ref{tab:learntropy-lr}). The
mechanism reveals a Piagetian inversion: the expert with the \emph{lowest}
learntropy ($\sim$2.4) and lowest LR ($1.8 \times 10^{-6}$, assimilation
mode) becomes the \emph{most critical} ($\Delta$PPL +39.7 when removed),
while the expert with the \emph{highest} learntropy ($\sim$5.0) and highest
LR ($2.4 \times 10^{-5}$, accommodation mode) becomes dormant ($\Delta$PPL
$<$0.1). Slow, stable learning crystallizes into indispensable knowledge;
over-accommodation prevents crystallization. The overall $M_{ij}$ CV reaches
1.25---higher than any previous experiment.

\begin{table}[h]
\centering
\caption{Learntropy-LR experiment: 4 experts with per-expert LR modulation
on Qwen3-8B. 3 of 4 experts active (vs.\ 2/4 in all previous experiments).}
\label{tab:learntropy-lr}
\begin{tabular}{@{}lcrrrrr@{}}
\toprule
Expert & Mode & LR & Medical & Code & Legal & Conv. \\
\midrule
E0 & Accomm. & $2.4 \times 10^{-5}$ & +0.06 & $-$0.02 & +0.00 & +0.14 \\
E1 & Moderate & $5.8 \times 10^{-6}$ & +8.51 & +11.67 & +11.21 & +21.19 \\
E2 & Assim. & $1.8 \times 10^{-6}$ & +27.19 & +38.96 & +31.85 & +70.01 \\
E3 & Moderate & $8.2 \times 10^{-6}$ & +2.58 & +3.30 & +2.68 & +5.28 \\
\bottomrule
\end{tabular}
\end{table}

\subsection{Path Forward: Learned Training Policies}

The three knobs---fork depth, fork timing, learning rate---define a compact
decision space. A training controller observes learntropy statistics (mean,
variance, kurtosis per expert per layer) and selects actions. Our 25+
experiment trajectories provide offline calibration data. The path proceeds
in three steps:

\emph{Step 1} (offline): extract learntropy statistics from existing
checkpoints and correlate with $M_{ij}$ outcomes to calibrate thresholds.

\emph{Step 2} (Bayesian optimization): fit a Gaussian Process surrogate over
the $\sim$15 controller parameters and optimize using short proxy runs (5K
steps) as evaluations.

\emph{Step 3} (learned policy): graduate from BO to a meta-learned controller
(Population-Based Training~\cite{jaderberg2017pbt} or bandit) that adapts
online during training.
A rule-based version is feasible now; a learned version
(Population-Based Training~\cite{jaderberg2017pbt} or bandit optimization)
would enable fully autonomous expert-tree growth.

% ============================================================================
\appendix

\section{Mathematical Details}
\label{app:math}

\subsection{Learntropy Bimodality Test}

Given an expert $e$ processing tokens $\{x_1, \ldots, x_T\}$, we compute the
per-token loss vector $\ell_e = [\ell_e(x_1), \ldots, \ell_e(x_T)]$ and apply
Hartigan's dip test~\cite{hartigan1985dip} for unimodality. The expert is a
split candidate when the dip statistic exceeds a threshold $\delta$:
\begin{equation}
    \text{DIP}(\ell_e) > \delta \implies \text{split expert } e
\end{equation}
We set $\delta = 0.05$ (the conventional significance level for rejecting
unimodality).

\subsection{Rank Growth Criterion}

After a split, each child adapter begins at rank $r = 1$. Let
$W_{\text{parent}}$ be the parent's weight matrix and $\Delta W = BA$ the LoRA
perturbation at current rank $r$. We compute the relative reconstruction error:
\begin{equation}
    \epsilon_r = \frac{\|W_{\text{target}} - W_{\text{parent}} - BA\|_F}{\|W_{\text{target}} - W_{\text{parent}}\|_F}
\end{equation}
where $W_{\text{target}}$ is the optimal weight matrix estimated from training
data. When $\epsilon_r > \epsilon_{\max}$, the rank is doubled: $r \leftarrow 2r$.

\subsection{Combined Training Objective}

The full training loss for each expert combines the standard language modeling
loss with the contrastive term:
\begin{equation}
    \mathcal{L} = \mathcal{L}_{\text{LM}} + \lambda_c \sum_{i < j} \max\bigl(0,\; \cos(\theta_{ij}) - \tau\bigr)
\end{equation}
We use $\lambda_c = 0.1$ and $\tau = 0.85$ in all experiments.

\section{Shannon's Entropy Argument for Zipf-Aligned Routing}
\label{app:shannon}

Shannon's analysis of English~\cite{shannon1948} established that natural
language has a specific entropy structure: highly redundant on average
($\sim$1 bit per character), but with information concentrated in sparse,
high-surprise moments. This is the source of Zipf's law---common tokens carry
little information and appear frequently, while rare tokens carry much
information and appear rarely.

Standard MoE load-balancing ignores this structure: it allocates equal capacity
to every expert regardless of information content. \tok{}'s argument is a
Shannon argument: parameter allocation should reflect information density. When
specialization depth follows Zipf's law, so does adapter rank, and therefore so
does storage footprint. The most frequently needed adapters are the smallest;
the largest adapters are needed least often. This access pattern is precisely
what tiered storage hierarchies exploit---not by coincidence, but because both
language and memory hierarchies are shaped by power-law economics.

At current scale (Qwen3-1.7B, adapters totaling $<$20\,MB), tiered storage is
unnecessary---everything fits in VRAM. Tiered deployment becomes relevant at
three thresholds: (1)~larger cores (70B+) where VRAM is contested,
(2)~thousands of community-contributed experts, and (3)~edge devices with
6--8\,GB total memory.

\section{Tiered Storage Validation}
\label{app:tiered}

We conducted 121 inference experiments on rented GPU instances (NVIDIA RTX PRO
6000 Blackwell, OLMoE-1B-7B) to calibrate a hardware-realistic cache simulator.

\begin{table}[h]
\centering
\caption{Memory tier economics (Netherlands retail, March 2026, incl.\ VAT).}
\label{tab:tiers}
\begin{tabular}{@{}lrrr@{}}
\toprule
Medium & Capacity & Bandwidth & Cost/GB \\
\midrule
GPU VRAM (GDDR7) & 48\,GB & $\sim$1.8\,TB/s & (fixed at mfg.) \\
DDR5-5600 & 256\,GB & $\sim$90\,GB/s & $\sim$\texteuro{}12/GB \\
NVMe Gen5 & 2--4\,TB & $\sim$14\,GB/s & $\sim$\texteuro{}0.09/GB \\
NVMe Gen4 & 2--4\,TB & $\sim$7\,GB/s & $\sim$\texteuro{}0.07/GB \\
\bottomrule
\end{tabular}
\end{table}

\begin{table}[h]
\centering
\caption{Expert load time from Gen5 NVMe at 14\,GB/s, vs.\ per-token compute
time at 134\,tok/s (7.5\,ms/token).}
\label{tab:load-times}
\begin{tabular}{@{}lrrl@{}}
\toprule
Adapter rank & Size & Load time & Critical path? \\
\midrule
Rank-1 & 4\,KB & 0.3\,$\mu$s & No (CPU cache) \\
Rank-16 & 64\,KB & 4.6\,$\mu$s & No ($<$0.1\% of token time) \\
Rank-64 & 256\,KB & 18\,$\mu$s & No ($<$0.3\% of token time) \\
Full expert (OLMoE) & 16\,MB & 1.1\,ms & Marginal (15\% of token time) \\
\bottomrule
\end{tabular}
\end{table}

Even at rank-64, adapter loads consume $<$0.3\% of a token's compute budget.
Tiered storage is not a compromise---it is invisible at \tok{} adapter sizes.

% ============================================================================
\begin{thebibliography}{99}
\bibitem{hu2022lora} E. Hu et al. LoRA: Low-Rank Adaptation of Large Language Models. \textit{ICLR}, 2022.
\bibitem{geva2021transformer} M. Geva, R. Schuster, J. Berant, O. Levy. Transformer Feed-Forward Layers Are Key-Value Memories. \textit{EMNLP}, 2021.
\bibitem{li2018intrinsic} C. Li et al. Measuring the Intrinsic Dimension of Objective Landscapes. \textit{ICLR}, 2018.
\bibitem{aghajanyan2021intrinsic} A. Aghajanyan et al. Intrinsic Dimensionality Explains the Effectiveness of Language Model Fine-Tuning. \textit{ACL}, 2021.
\bibitem{komatsuzaki2023sparse} A. Komatsuzaki et al. Sparse Upcycling. \textit{arXiv:2212.05055}, 2023.
\bibitem{bena2025modularity} G. B\'{e}na, D. Goodman. Dynamics of Specialization Under Resource Constraints. \textit{Nature Communications}, 2025.
\bibitem{l2r2026} L2R: Low-Rank Lipschitz-Controlled Routing. \textit{arXiv:2601.21349}, 2026.
\bibitem{shannon1948} C. E. Shannon. A Mathematical Theory of Communication. \textit{Bell System Technical Journal}, 27(3):379--423, 1948.
\bibitem{hartigan1985dip} J. A. Hartigan, P. M. Hartigan. The Dip Test of Unimodality. \textit{Annals of Statistics}, 13(1):70--84, 1985.
\bibitem{piaget1952} J. Piaget. \textit{The Origins of Intelligence in Children}. International Universities Press, 1952.
\bibitem{vygotsky1978} L. S. Vygotsky. \textit{Mind in Society: The Development of Higher Psychological Processes}. Harvard University Press, 1978.
\bibitem{wozniak2018learntropy} P. A. Wozniak. Learntropy. \textit{SuperMemo Guru}, 2018. \url{https://supermemo.guru/wiki/Learntropy}.
\bibitem{standing2026} The Illusion of Specialization: Standing Committee in Mixture-of-Experts. \textit{arXiv:2601.03425}, 2026.
\bibitem{deepseek2024moe} X. Dai et al. DeepSeekMoE: Towards Ultimate Expert Specialization. \textit{ACL}, 2024.
\bibitem{deepseekv3} DeepSeek-AI. DeepSeek-V3 Technical Report. \textit{arXiv:2412.19437}, 2024.
\bibitem{hydralora2024} T. Tian et al. HydraLoRA: An Asymmetric LoRA Architecture for Efficient Fine-Tuning. \textit{NeurIPS}, 2024.
\bibitem{molora2026} MoLoRA: Composable Specialization via Per-Token Adapter Routing. \textit{arXiv:2603.15965}, 2026.
\bibitem{basicrefinement2025} Decoding Knowledge Attribution in MoE: A Framework of Basic-Refinement Collaboration. \textit{ACL}, 2025.
\bibitem{openmoe2024} F. Xue et al. OpenMoE: An Early Effort on Open Mixture-of-Experts Language Models. \textit{ICML}, 2024.
\bibitem{mixtral2024} A. Jiang et al. Mixtral of Experts. \textit{arXiv:2401.04088}, 2024.
\bibitem{domaindriver2026} What Gets Activated: Uncovering Domain and Driver Experts in MoE Language Models. \textit{arXiv:2601.10159}, 2026.
\bibitem{jang2017gumbel} E. Jang, S. Gu, B. Poole. Categorical Reparameterization with Gumbel-Softmax. \textit{ICLR}, 2017.
\bibitem{monet2025} S. Park et al. Monet: Mixture of Monosemantic Experts for Transformers. \textit{ICLR}, 2025.
\bibitem{idamoe2025} IDA-MoE: Input Domain Aware Mixture-of-Experts. \textit{arXiv:2510.16448}, 2025.
\bibitem{jaderberg2017pbt} M. Jaderberg et al. Population Based Training of Neural Networks. \textit{arXiv:1711.09846}, 2017.
\end{thebibliography}

\end{document}
